\documentclass[UTF8,zihao=-4]{ctexart}

\usepackage[a4paper,margin=2.05cm]{geometry}
\usepackage{xcolor}
\usepackage{enumitem}
\usepackage{hyperref}
\usepackage{longtable}
\usepackage{array}
\usepackage{booktabs}
\usepackage{listings}
\usepackage{parskip}

\definecolor{TitleBlue}{HTML}{1F4E79}
\definecolor{AnswerGray}{HTML}{333333}
\definecolor{SoftGray}{HTML}{F4F6F8}

\hypersetup{
  colorlinks=true,
  linkcolor=TitleBlue,
  urlcolor=TitleBlue
}

\setlist[itemize]{leftmargin=2em,itemsep=0.25em,topsep=0.25em}
\setlist[enumerate]{leftmargin=2.2em,itemsep=0.25em,topsep=0.25em}
\setlength{\parindent}{0pt}
\renewcommand{\baselinestretch}{1.12}

\lstset{
  basicstyle=\ttfamily\small,
  breaklines=true,
  columns=fullflexible,
  frame=single,
  backgroundcolor=\color{SoftGray},
  xleftmargin=0.5em,
  xrightmargin=0.5em
}

\newcounter{qcounter}
\newcommand{\chaptertitle}[1]{%
  \section*{\color{TitleBlue}#1}
  \addcontentsline{toc}{section}{#1}
}
\newcommand{\question}[1]{%
  \refstepcounter{qcounter}
  \subsection*{题 \theqcounter：#1}
  \addcontentsline{toc}{subsection}{题 \theqcounter：#1}
}
\newenvironment{answer}{%
  \textbf{答案：}\par\small\color{AnswerGray}
}{%
  \par\normalsize\color{black}
}

\title{\textbf{软件工程与计算 II 复习提纲含答案版}}
\author{}
\date{}

\begin{document}
\pagestyle{plain}
\maketitle
\tableofcontents
\newpage

\section*{\color{TitleBlue}整理说明}
\addcontentsline{toc}{section}{整理说明}
本文件依据工作区根目录的《复习提纲.pdf》整理。原提纲是 26 页课件式题目清单，没有给出标准答案。本版按原 PDF 的章节顺序补充作答要点、答题模板和典型例题答案。

原 PDF 中出现“结合实验”的题目，但未提供具体实验项目名、团队分工、仓库记录、测试报告和缺陷记录。因此本版给出可直接套用的实验答题模板，不编造具体项目事实；考试或复习时把模板中的项目名称、模块名称、工具链和实际过程替换为自己的实验记录。

\chaptertitle{第一、二章：软件工程基础}

\question{名词解释：软件工程}
\begin{answer}
软件工程是把系统化、规范化、可量化的方法应用于软件开发、运行和维护，并研究这些方法的学科。

展开时写三层：
\begin{enumerate}
  \item 方法层：系统化、规范化、可量化。
  \item 对象层：软件制品包括程序、数据、文档和知识，不只是代码。
  \item 目标层：在多人协作和长期演化中控制质量、成本、进度和风险，交付可维护的软件系统。
\end{enumerate}
\end{answer}

\question{简答：1950s 到 2000s 软件工程发展的特点}
\begin{answer}
按时间线作答：
\begin{longtable}{p{0.16\linewidth}p{0.75\linewidth}}
\toprule
阶段 & 特点 \\
\midrule
1950s & 商业计算机和早期高级语言出现，软件主要依附硬件，程序规模小，开发依赖个人经验。\\
1960s & 软件规模和复杂度上升，延期、超支、质量低和维护困难集中暴露，1968 年 NATO 会议提出“软件危机”和“软件工程”。\\
1970s & 强调工程化过程，瀑布模型、结构化分析和结构化设计兴起，文档和阶段评审成为控制手段。\\
1980s & 软件生命周期、质量管理、配置管理和成熟度模型发展，团队协作和过程改进成为重点。\\
1990s & 面向对象方法、UML、构件复用和迭代思想普及，软件系统更强调可扩展和可维护。\\
2000s & 敏捷方法、持续集成、开源协作、Web 系统和服务化发展，反馈速度、自动化测试和持续交付成为重要工程能力。\\
\bottomrule
\end{longtable}
收束句：这条发展线的核心是从“个人写程序”转向“团队用工程方法管理复杂度、变化和质量”。
\end{answer}

\chaptertitle{第四章：项目启动与项目管理}

\question{如何管理团队？在实验中采取哪些办法？有哪些经验？}
\begin{answer}
团队管理从目标、角色、计划、沟通、质量和风险六方面写：
\begin{enumerate}
  \item 明确共同目标：把项目目标拆成迭代目标、里程碑和可交付物。
  \item 明确角色分工：项目负责人、需求负责人、架构或设计负责人、开发负责人、测试负责人、配置管理员。
  \item 制定计划：任务拆分、估算工作量、排定优先级和截止时间。
  \item 建立沟通机制：例会、看板、Issue、Pull Request 和评审记录。
  \item 控制质量：需求评审、设计评审、代码审查、单元测试、集成测试和回归测试。
  \item 管理风险：识别关键成员缺席、需求变更、技术不熟、接口延期等风险，提前准备应对措施。
\end{enumerate}

结合实验的答题模板：
\begin{quote}
在实验项目中，我们将任务拆为需求、设计、编码、测试和文档五类；通过 Git 仓库、Issue、分支和 Pull Request 跟踪任务；每次迭代结束进行需求和代码评审；对核心功能编写单元测试和集成测试；对延期或接口变化风险设置负责人和截止时间。经验是：分工要落到具体交付物，沟通要有记录，合并代码前要评审和测试。
\end{quote}
\end{answer}

\question{团队结构有哪几种？}
\begin{answer}
可按控制方式和沟通方式分类：
\begin{itemize}
  \item 主程序员式团队：核心成员负责总体设计和关键决策，适合小团队和技术集中项目。
  \item 民主式团队：成员平等讨论，适合探索性强、创新性强的任务。
  \item 分层式团队：项目经理、模块负责人、开发和测试分层协作，适合规模较大的项目。
  \item 敏捷小队：跨职能成员共同负责迭代交付，强调频繁反馈和持续改进。
\end{itemize}
答题时说明：团队结构要服务于项目规模、成员能力、需求稳定性和沟通成本。
\end{answer}

\question{质量保障有哪些措施？结合实验说明}
\begin{answer}
软件质量保证 SQA 关注过程和产品是否满足质量要求。常见措施：
\begin{itemize}
  \item 需求阶段：需求评审、术语统一、可验证性检查。
  \item 设计阶段：架构评审、接口评审、设计类图和顺序图检查。
  \item 构造阶段：代码审查、静态分析、单元测试、持续集成。
  \item 测试阶段：功能测试、集成测试、回归测试、缺陷跟踪。
  \item 发布阶段：配置审计、版本核对、冒烟测试和回滚预案。
\end{itemize}

结合实验的答题模板：
\begin{quote}
在实验项目中，质量保障贯穿需求、设计、编码、测试和发布。需求文档完成后进行组内评审；设计阶段检查分层结构和接口依赖；编码阶段使用分支开发、Pull Request 和代码审查；核心业务编写单元测试，跨模块流程编写集成测试；每次发布前核对版本、执行冒烟测试并记录缺陷。
\end{quote}
\end{answer}

\question{配置管理有哪些活动？实验中如何进行配置管理？}
\begin{answer}
软件配置管理 SCM 的目标是识别配置项、控制变更、记录状态、审计配置并保证产品与规格一致。

活动包括：
\begin{enumerate}
  \item 配置项识别：需求文档、设计文档、源代码、测试脚本、部署脚本、数据库脚本、配置文件和发布包。
  \item 版本管理：分支、标签、提交记录和版本号。
  \item 变更控制：提交变更请求，影响分析，审批，执行，验证。
  \item 配置审计：检查代码、配置、脚本、依赖和部署清单是否一致。
  \item 状态报告：记录每个配置项的版本、负责人、状态和变更历史。
  \item 发布管理：建立发布基线，执行发布检查和回滚预案。
\end{enumerate}

实验答题模板：
\begin{quote}
我们将需求文档、设计文档、源代码、测试用例、数据库脚本和部署配置纳入 Git 管理；主分支保持可运行，功能在独立分支开发；合并前通过 Pull Request 审查；重要版本打 Tag；需求或接口变化通过 Issue 记录原因、影响范围和处理结果；发布前核对配置和版本并执行冒烟测试。
\end{quote}
\end{answer}

\question{示例：结合实验说明一个项目的质量保障包括哪些活动}
\begin{answer}
标准结构：
\begin{enumerate}
  \item 需求质量保障：需求评审、冲突检查、可验证性检查。
  \item 设计质量保障：架构评审、接口评审、职责分配检查。
  \item 代码质量保障：编码规范、代码审查、静态分析、重构。
  \item 测试质量保障：单元测试、集成测试、系统测试、回归测试。
  \item 配置质量保障：版本管理、变更控制、配置审计、状态报告。
  \item 发布质量保障：发布基线、部署清单、冒烟测试、回滚计划。
\end{enumerate}
答题句：质量保障不是只做测试，而是通过评审、测试、度量、配置管理和过程审计共同降低缺陷进入后续阶段的风险。
\end{answer}

\chaptertitle{第五章：需求工程}

\question{名词解释：需求}
\begin{answer}
需求是利益相关者为解决问题或达到目标而要求系统具备的能力、行为、约束或质量属性。需求既包括系统要做什么，也包括系统必须满足的性能、安全、可靠性、接口、数据和业务规则约束。
\end{answer}

\question{区分需求的三个层次，并给出实例}
\begin{answer}
三个层次：
\begin{longtable}{p{0.2\linewidth}p{0.33\linewidth}p{0.36\linewidth}}
\toprule
层次 & 回答的问题 & 校园卡充值例子 \\
\midrule
业务需求 & 为什么建设系统 & 提高校园卡充值自助化比例，降低人工窗口压力。\\
用户需求 & 用户要完成什么任务 & 学生希望在线给校园卡充值并查看充值结果。\\
系统级需求 & 系统在条件下必须做什么 & 当支付平台返回支付成功结果时，系统应更新校园卡余额并生成充值记录。\\
\bottomrule
\end{longtable}

判定方法：看需求回答的是“为什么建系统”“用户想完成什么任务”还是“系统必须提供什么可验证行为”。
\end{answer}

\question{掌握需求类型，并对实例判定类型}
\begin{answer}
常见类型：
\begin{itemize}
  \item 功能需求：系统要做什么，如“系统应支持用户查询余额”。
  \item 性能需求：速度、容量、吞吐量、响应时间，如“查询响应时间不超过 2 秒”。
  \item 质量属性：安全性、可靠性、可维护性、易用性、可用性。
  \item 对外接口需求：与外部系统、设备、用户界面的交互方式，如“通过 HTTPS 调用支付平台接口”。
  \item 约束：技术栈、法规、部署环境、开发过程限制。
  \item 数据需求：数据字段、格式、关系、完整性规则和保留期限。
\end{itemize}

判定句式：
\begin{quote}
该需求属于功能需求，因为它描述系统在某个输入或事件下必须执行的业务行为。该需求属于数据需求，因为它规定了字段格式、取值范围或持久化内容。该需求属于对外接口需求，因为它规定了本系统与外部系统的调用方式、数据格式或异常处理。
\end{quote}
\end{answer}

\chaptertitle{第六章：需求分析建模}

\question{给定描述时，如何建立用例图、分析类图、系统顺序图和状态图？}
\begin{answer}
建模顺序：
\begin{enumerate}
  \item 用例图：识别参与者、系统边界、用户目标、include 和 extend 关系。
  \item 分析类图：从需求文本中提取领域名词，保留系统需要长期记录或参与业务规则的概念；只写属性和关联，不写方法。
  \item 系统顺序图：把系统看成黑盒，画参与者与系统之间的事件、返回和异常流程，不画内部对象。
  \item 状态图：选择生命周期明显的对象，如订单、预约、借阅、校园卡；列出状态、事件、条件和动作。
\end{enumerate}
检查点：边界清楚、粒度适中、概念类不混入界面控件或数据库表、顺序图不画内部实现。
\end{answer}

\question{图书馆系统：建立用例模型并说明建模过程}
\begin{answer}
建模过程：
\begin{enumerate}
  \item 识别参与者：借书人、管理员、图书馆管理者、外部大学数据库。借书人有教职工借书人、研究生借书人、本科生借书人三个特化。
  \item 识别系统边界：简化大学图书馆系统。
  \item 提取用户目标：检索书名、查询可借状态、借书、还书、预约、跟踪新书、打印标题表、检查逾期未还、下载借书人信息。
  \item 组织关系：借书人参与检索、借书、预约；管理员参与借书、还书、跟踪新书；管理者参与打印报表和检查逾期；系统访问外部大学数据库下载借书人信息。
\end{enumerate}

文字版用例模型：
\begin{longtable}{p{0.22\linewidth}p{0.68\linewidth}}
\toprule
参与者 & 用例 \\
\midrule
借书人 & 检索书名、查询图书可借状态、预约图书、办理借书。\\
管理员 & 办理借书、办理还书、跟踪新书到达。\\
图书馆管理者 & 打印图书标题表、在线检查过期未还图书。\\
外部大学数据库 & 提供借书人信息下载服务。\\
\bottomrule
\end{longtable}

include 和 extend 的写法：办理借书包含“登记借出记录”和“关联图书副本”；当所有副本已借出时，“预约图书”作为条件性扩展发生。
\end{answer}

\question{ATM 系统：建立领域类图，即分析类图}
\begin{answer}
领域概念和属性：
\begin{longtable}{p{0.24\linewidth}p{0.62\linewidth}}
\toprule
概念类 & 属性或说明 \\
\midrule
ATM & 设备编号、位置、状态。\\
显示屏幕 & 显示内容。\\
输入键盘 & 数字键、特殊符号键。\\
银行卡读卡器 & 读卡状态。\\
存款插槽 & 接收存款。\\
收据打印机 & 打印收据。\\
客户 & 客户编号、姓名、安全级别。\\
银行卡 & 卡号、状态。\\
账户 & 账号、余额、账户状态。\\
事务 & 事务编号、时间、类型、金额、结果。\\
存款事务 & 存款金额。\\
取款事务 & 取款金额。\\
余额查询事务 & 查询时间。\\
安全系统 & PIN 验证规则、安全级别分配。\\
账户系统接口 & 账户更新、余额查询。\\
\bottomrule
\end{longtable}

关系：
\begin{itemize}
  \item ATM 由显示屏幕、输入键盘、读卡器、存款插槽、收据打印机组成。
  \item 客户持有银行卡，银行卡关联一个或多个账户。
  \item 客户通过 ATM 发起事务，事务分为存款、取款、余额查询。
  \item 每次事务执行前调用安全系统验证 PIN。
  \item 存款、取款和余额查询通过账户系统接口访问或更新账户。
\end{itemize}
分析类图只表示问题域概念、属性和关联，不写 Controller、Service、Repository 或数据库表。
\end{answer}

\chaptertitle{第七章：需求规格说明与测试}

\question{为什么需要需求规格说明？结合实验说明}
\begin{answer}
需求规格说明 SRS 的作用：
\begin{enumerate}
  \item 建立共同理解：让用户、开发、测试和管理者对系统能力达成一致。
  \item 支持设计和实现：为架构、接口、数据和测试提供输入。
  \item 支持验证：每条需求都应可测试、可追踪。
  \item 控制变更：变更发生时能分析影响范围。
  \item 形成合同或验收依据：明确系统交付边界。
\end{enumerate}

实验答题模板：
\begin{quote}
在实验项目中，SRS 记录了功能需求、数据需求、接口需求、非功能需求和验收标准。设计阶段依据 SRS 划分模块和接口；测试阶段依据 SRS 设计功能测试用例；需求变更时根据编号追踪到设计、代码和测试，避免遗漏。
\end{quote}
\end{answer}

\question{对给定需求规格说明片段，如何判定并修正错误？}
\begin{answer}
检查维度：
\begin{itemize}
  \item 歧义：出现“尽快”“方便”“大量”“正常情况”等不可验证词。
  \item 不完整：缺少触发条件、输入、输出、异常流程或边界值。
  \item 不一致：同一术语、状态或业务规则前后冲突。
  \item 不可验证：没有明确阈值、条件和期望结果。
  \item 混入设计实现：把数据库表、按钮布局、类名等实现细节写成需求。
  \item 多需求混写：一个条目同时描述多个行为。
\end{itemize}

修正模板：
\begin{quote}
当 \texttt{<触发条件>} 发生，且 \texttt{<前置条件>} 成立时，系统应 \texttt{<可观察响应>}；若 \texttt{<异常条件>} 发生，系统应 \texttt{<异常响应>}。
\end{quote}

例子：把“系统要很快完成查询”改为“当用户提交余额查询请求时，系统应在 2 秒内返回当前余额或失败原因，不包含外部网络故障时间”。
\end{answer}

\question{对给定需求示例，如何设计功能测试用例？}
\begin{answer}
功能测试基于规格，不看内部代码。设计方法：
\begin{enumerate}
  \item 等价类划分：有效输入类和无效输入类。
  \item 边界值分析：最小值、最大值、刚好越界值。
  \item 决策表：多个条件组合导致不同结果。
  \item 状态转换：对象在不同状态下接受事件产生不同响应。
  \item 场景测试：覆盖主成功流程和关键异常流程。
\end{enumerate}

测试用例表：
\begin{longtable}{p{0.14\linewidth}p{0.26\linewidth}p{0.24\linewidth}p{0.26\linewidth}}
\toprule
编号 & 输入/前置条件 & 覆盖点 & 预期结果 \\
\midrule
TC-01 & 正常用户、有效金额、支付成功 & 主成功流程 & 更新余额，生成成功记录。\\
TC-02 & 金额为 0 & 无效等价类 & 拒绝请求并提示金额非法。\\
TC-03 & 校园卡已挂失 & 状态异常 & 拒绝充值，不更新余额。\\
TC-04 & 支付平台返回失败 & 外部接口异常 & 记录失败原因并提示用户。\\
\bottomrule
\end{longtable}
\end{answer}

\chaptertitle{第八章：软件设计基础}

\question{名词解释：软件设计}
\begin{answer}
软件设计是把需求转换为可实现的软件解决方案的过程。它确定系统的体系结构、模块划分、接口、数据结构、类和对象协作、算法、异常处理和部署约束。设计回答“如何实现需求”。
\end{answer}

\question{软件设计的核心思想是什么？}
\begin{answer}
核心思想是管理复杂度和变化。主要手段：
\begin{itemize}
  \item 分解：把复杂系统拆成模块、层、组件和类。
  \item 抽象：保留关键概念，隐藏不必要细节。
  \item 信息隐藏：隐藏重要设计决策和变化点。
  \item 模块化：通过高内聚、低耦合降低修改传播。
  \item 接口化：依赖稳定抽象而不是具体实现。
\end{itemize}
\end{answer}

\question{软件工程设计有哪三个层次？各层主要思想是什么？}
\begin{answer}
三个层次：
\begin{longtable}{p{0.22\linewidth}p{0.67\linewidth}}
\toprule
层次 & 主要思想 \\
\midrule
体系结构设计 & 确定系统整体结构、分层、组件、连接件、部署节点和技术约束。\\
详细设计 & 确定类、接口、职责分配、对象协作、控制风格和方法级结构。\\
代码设计 & 关注变量、语句、控制结构、异常处理、命名、布局、可读性和可维护性。\\
\bottomrule
\end{longtable}
\end{answer}

\chaptertitle{第九、十章：体系结构}

\question{体系结构的概念}
\begin{answer}
软件体系结构是系统的高层结构，包含组件、连接件和配置。组件承担计算、存储或业务职责；连接件定义组件之间如何交互，如方法调用、REST API、消息队列；配置说明组件和连接件如何组合成系统。

答题句：体系结构决策影响系统的可维护性、性能、安全性、可扩展性和部署方式，一旦落地修改成本高。
\end{answer}

\question{体系结构风格的优缺点}
\begin{answer}
常用风格：
\begin{longtable}{p{0.18\linewidth}p{0.36\linewidth}p{0.36\linewidth}}
\toprule
风格 & 优点 & 缺点 \\
\midrule
分层风格 & 职责清晰，便于维护、测试和替换下层实现。 & 层次过多会增加调用开销，错误依赖会破坏结构。\\
MVC & 分离模型、视图和控制，便于界面变化和业务复用。 & 控制器容易膨胀，前后端职责边界需要明确。\\
管道-过滤器 & 适合数据流处理，过滤器可复用。 & 不适合强交互和共享状态复杂的系统。\\
客户端-服务器 & 部署边界清楚，客户端和服务器职责分离。 & 网络延迟、服务端瓶颈和安全问题需要控制。\\
微服务 & 服务独立部署和扩展，适合大型组织。 & 分布式事务、监控、运维和接口治理复杂。\\
\bottomrule
\end{longtable}
\end{answer}

\question{体系结构设计过程}
\begin{answer}
作答步骤：
\begin{enumerate}
  \item 分析需求和质量属性。
  \item 识别关键场景和架构风险。
  \item 选择体系结构风格。
  \item 划分逻辑层、组件和包。
  \item 设计层间接口和依赖方向。
  \item 设计物理部署和通信方式。
  \item 用核心用例、集成测试和评审验证架构。
\end{enumerate}
\end{answer}

\question{包的原则}
\begin{answer}
包用于组织类和接口，原则包括：
\begin{itemize}
  \item 单一职责：同一包内元素围绕同一职责。
  \item 稳定依赖：不稳定包依赖稳定抽象，避免稳定包依赖易变实现。
  \item 无环依赖：包之间不形成循环依赖。
  \item 依赖接口：上层依赖下层接口而不是实现类。
  \item 按层和领域组织：Controller、Service、Repository、Domain 等职责边界清楚。
\end{itemize}
\end{answer}

\question{体系结构构件之间接口的定义}
\begin{answer}
接口定义组件之间的服务契约，至少说明：
\begin{itemize}
  \item 接口所在包名和接口名。
  \item 操作名、参数、返回值和异常。
  \item 前置条件、后置条件和错误处理。
  \item 调用方向和依赖方向。
\end{itemize}
分层系统中，Controller 依赖 Service 接口，ServiceImpl 依赖 Repository 接口，Repository 不处理业务规则。
\end{answer}

\question{体系结构开发集成测试用例：Stub 和 Driver}
\begin{answer}
集成测试检查模块协作和接口契约。
\begin{itemize}
  \item Stub：被测模块调用的下层替身。上层已完成、下层未完成时，用 Stub 模拟下层返回。
  \item Driver：调用被测模块的上层替身。下层已完成、上层未完成时，用 Driver 触发被测模块。
\end{itemize}

例子：测试 Service 时，Repository 未完成，使用 Repository Stub 返回固定数据；测试 Repository 时，Controller 未完成，使用 Driver 调用 Repository 方法并检查数据库结果。
\end{answer}

\question{薪酬管理系统 Web 分层物理包图如何作答？}
\begin{answer}
物理包图必须体现客户端、网络通信、服务器端分层和功能模块。

文字版结构：
\begin{itemize}
  \item 客户端浏览器：\texttt{html}、\texttt{css}、\texttt{js}。
  \item 网络通信：HTTP，基于 REST 风格 API。
  \item 服务器表示入口：\texttt{controller} 包，包含员工、薪资和社会保险录入等 Controller。
  \item 逻辑层：\texttt{service} 包和 \texttt{service.impl} 包，如员工管理、薪资计算、社会保险录入。
  \item 数据层：\texttt{repository} 包和 \texttt{repository.impl} 包。
  \item 领域对象：\texttt{domain} 包，如 \texttt{Employee}、\texttt{SalaryRecord}、\texttt{SocialInsuranceRecord}。
  \item 数据库：员工表、薪资表、社会保险表。
\end{itemize}

得分点：浏览器在客户端，逻辑层和数据层在服务器端；客户端与服务器通过 HTTP/REST 通信；服务器内部体现 Controller、Service、Repository 和数据库依赖。
\end{answer}

\question{“社会保险录入”页面的展示层、逻辑层、数据层接口声明}
\begin{answer}
展示层与逻辑层之间的接口放在 \texttt{payroll.service} 包：
\begin{lstlisting}[language=Java]
package payroll.service;

public interface SocialInsuranceService {
    SocialInsuranceResult createRecord(SocialInsuranceCommand command);
    SocialInsuranceRecordDTO findByEmployeeId(String employeeId);
}
\end{lstlisting}

逻辑层与数据层之间的接口放在 \texttt{payroll.repository} 包：
\begin{lstlisting}[language=Java]
package payroll.repository;

public interface SocialInsuranceRepository {
    void save(SocialInsuranceRecord record);
    SocialInsuranceRecord findByEmployeeId(String employeeId);
}
\end{lstlisting}

如需员工信息校验，再写员工仓储接口：
\begin{lstlisting}[language=Java]
package payroll.repository;

public interface EmployeeRepository {
    Employee findById(String employeeId);
}
\end{lstlisting}

解释：Controller 只调用 \texttt{SocialInsuranceService}；ServiceImpl 校验员工和保险录入规则后调用 Repository；Repository 只负责数据访问，不写业务规则。
\end{answer}

\chaptertitle{第十一章：人机交互}

\question{名词解释：可用性}
\begin{answer}
可用性表示特定用户在特定使用环境下，使用系统完成特定任务时达到有效性、效率和满意度的程度。课程答题常写五项属性：易学性、效率、易记性、错误率与恢复、满意度。
\end{answer}

\question{列出至少 5 个界面设计注意事项并解释}
\begin{answer}
可写以下原则：
\begin{enumerate}
  \item 一致性：相同功能在不同页面使用一致的术语、位置、颜色和交互方式。
  \item 反馈：用户操作后系统立即显示状态、结果或进度。
  \item 识别优于回忆：让可选项、历史记录和状态可见，减少用户记忆负担。
  \item 错误预防：通过约束输入、默认值、确认框和撤销减少错误。
  \item 用户控制和自由：支持取消、返回、撤销和修改。
  \item Fitts 定律：高频操作按钮要足够大并靠近当前任务区域。
  \item 任务匹配：界面按用户任务组织，而不是按数据库表或技术模块组织。
\end{enumerate}
\end{answer}

\question{分析图中两个界面体现或违反的人机交互设计原则}
\begin{answer}
左图主要问题：
\begin{itemize}
  \item 违反识别优于回忆：用户需要通过滚动条逐个查找样式，列表不可见。
  \item 反馈不足：当前选择和最终效果之间的关系不够直接。
  \item 效率低：用户需要反复滚动才能比较样式。
  \item 控件匹配差：横向滚动条不适合浏览大量样式名称。
\end{itemize}

右图改进点：
\begin{itemize}
  \item 支持识别优于回忆：左侧列表直接显示多个样式名称。
  \item 反馈更明确：右侧预览随选择变化，用户能看到结果。
  \item 效率更高：列表和预览并列，减少来回操作。
  \item 仍需注意：若列表过长，应提供搜索、分类或更大的预览区域。
\end{itemize}
\end{answer}

\question{精神模型、差异性、导航、反馈、协作式设计}
\begin{answer}
答题要点：
\begin{itemize}
  \item 精神模型：用户对系统如何工作的理解。界面应符合用户的任务习惯和领域概念。
  \item 差异性：不同用户能力、经验、设备和使用场景不同，界面要兼顾新手、中级用户和专家用户。
  \item 导航：让用户知道自己在哪里、能去哪、如何返回和下一步是什么。
  \item 反馈：操作后及时显示状态、结果、错误原因或进度。
  \item 协作式设计：让用户、设计者、开发者共同参与需求澄清、原型评审和可用性测试。
\end{itemize}
\end{answer}

\chaptertitle{第十二章：详细设计}

\question{详细设计的出发点是什么？}
\begin{answer}
详细设计从分析模型、系统顺序图、架构约束和设计因素出发，把“系统必须做什么”转换为“对象如何协作完成”。重点是职责分配、对象协作、控制风格、类和接口设计、方法签名和测试设计。
\end{answer}

\question{职责分配、协作、控制风格分别是什么？}
\begin{answer}
\begin{itemize}
  \item 职责分配：决定哪个类负责知道信息、执行操作、创建对象或协调流程。
  \item 协作：对象之间通过消息调用共同完成用例。
  \item 控制风格：用例流程由谁主导。
\end{itemize}

控制风格：
\begin{longtable}{p{0.16\linewidth}p{0.32\linewidth}p{0.38\linewidth}}
\toprule
风格 & 特征 & 评价 \\
\midrule
集中式 & 一个控制对象掌握大量流程。 & 流程集中但控制器容易臃肿。\\
委托式 & 控制器保留主流程，子任务委托给专业对象。 & 职责清晰，常作为考试推荐结构。\\
分散式 & 流程散落在多个对象中。 & 单个对象小，但控制流难追踪。\\
\bottomrule
\end{longtable}
\end{answer}

\question{给定分析类图、系统顺序图和设计因素，如何建立设计类图或详细顺序图？}
\begin{answer}
步骤：
\begin{enumerate}
  \item 从系统顺序图中找系统事件。
  \item 为系统事件选择 Controller。
  \item 根据 GRASP 信息专家原则，把业务判断分给拥有数据的领域对象。
  \item 为外部依赖设计接口，如 Repository、Gateway、Service。
  \item 在设计类图中补充方法、可见性、接口、依赖和实现关系。
  \item 在详细顺序图中画 Controller、Service、Domain、Repository 的调用顺序。
  \item 用低耦合、高内聚、信息隐藏和可测试性检查设计。
\end{enumerate}
\end{answer}

\question{协作测试中的 MockObject}
\begin{answer}
MockObject 用来替代外部依赖，使测试集中在当前对象的协作行为上。它能验证：
\begin{itemize}
  \item 被测对象是否调用了正确依赖。
  \item 调用参数是否正确。
  \item 外部依赖返回成功或失败时，被测对象是否正确处理。
  \item 不访问真实数据库、支付平台或网络服务，从而让单元测试稳定、快速、可重复。
\end{itemize}
\end{answer}

\question{示例：根据 CD 推荐用例画设计类图和集中式控制顺序图}
\begin{answer}
原图给出顾客向系统提交查询请求，系统返回 CD 推荐列表，顾客选择 CD 后查询详细信息并购买 CD。集中式控制写法如下：
\begin{itemize}
  \item 设计类：\texttt{RecommendationController} 作为控制类；\texttt{Customer}、\texttt{Request}、\texttt{CDRecommendationList}、\texttt{CD} 作为领域类；\texttt{CDRepository} 作为数据访问接口。
  \item 关键方法：
  \begin{itemize}
    \item \texttt{queryRecommendations}
    \item \texttt{getDetail}
    \item \texttt{purchase}
  \end{itemize}
  \item 集中式顺序：Customer -> Controller -> Repository -> Controller -> Customer。后续“选定 CD”“查询详细信息”“购买 CD”也由 Controller 统一协调。
\end{itemize}

评价：集中式控制容易画、流程清楚，但 Controller 会积累推荐、查询和购买逻辑。改进时把推荐规则委托给 \texttt{RecommendationService}，把购买规则委托给 \texttt{OrderService}。
\end{answer}

\chaptertitle{第十三章：耦合、内聚与信息隐藏}

\question{名词解释：耦合与内聚}
\begin{answer}
耦合表示模块之间依赖强度。耦合越低，一个模块变化越不容易影响其他模块。

内聚表示模块内部职责相关程度。内聚越高，模块越围绕单一目标组织，越容易理解、测试和修改。
\end{answer}

\question{如何说明例子中的耦合程度与内聚，并给理由？}
\begin{answer}
耦合判断：
\begin{itemize}
  \item 内容耦合：直接访问或修改其他模块内部。
  \item 公共耦合：多个模块共享全局数据。
  \item 控制耦合：传递控制标志让被调用者选择行为。
  \item 印记耦合：传递整个对象或结构，但只使用其中一部分。
  \item 数据耦合：只通过必要参数交互。
\end{itemize}

内聚判断：
\begin{itemize}
  \item 偶然内聚：职责无明确关联。
  \item 逻辑内聚：用类型参数选择不同逻辑。
  \item 时间内聚：因同一时间执行而放在一起。
  \item 过程内聚：按流程顺序放在一起。
  \item 通信内聚：操作同一数据集合。
  \item 功能内聚：只完成一个明确功能。
  \item 信息内聚：围绕同一数据结构组织多个操作。
\end{itemize}
答题先给类型，再说明依赖方式或职责关系，再给修改影响。
\end{answer}

\question{信息隐藏的基本思想和两种常见信息隐藏决策}
\begin{answer}
信息隐藏的基本思想：模块隐藏一个重要设计决策，只暴露稳定接口。信息隐藏不是简单把字段设为 private，而是隐藏容易变化的决策。

常见决策：
\begin{enumerate}
  \item 数据结构或存储方式：用 Repository 接口隐藏数据库表、文件、缓存或远程服务。
  \item 算法或业务规则：用 Strategy、Service 或领域对象隐藏不同规则实现。
\end{enumerate}

判断信息隐藏好坏：外部模块是否依赖内部字段、具体类、算法步骤、数据库结构或第三方接口细节。依赖越少，信息隐藏越好。
\end{answer}

\chaptertitle{第十四章：模块化原则}

\question{Principles from Modularization 九条原则及修正方式}
\begin{answer}
\begin{longtable}{p{0.32\linewidth}p{0.54\linewidth}}
\toprule
原则 & 答题要点 \\
\midrule
Global Variables Consider Harmful & 避免全局变量和共享可变状态；改为参数、对象字段、依赖注入或配置对象。\\
To be Explicit & 依赖、输入、输出和异常要显式表达；不要依赖隐式全局状态。\\
Do not Repeat & 不重复知识和逻辑；重复代码提取方法、类或配置表。\\
Programming to Interface & 面向接口编程，依赖抽象而不是实现。\\
Law of Demeter & 只和直接朋友通信，避免长消息链。\\
Interface Segregation Principle & 接口小而专用，客户端不依赖自己不用的方法。\\
Liskov Substitution Principle & 子类必须能替换父类，不破坏父类语义。\\
Favor Composition Over Inheritance & 优先组合和委托，减少白盒复用风险。\\
Single Responsibility Principle & 一个模块只有一个变化理由。\\
\bottomrule
\end{longtable}

对给定代码题的作答结构：指出违反的原则，说明风险，给出改写结构，再补充测试。
\end{answer}

\chaptertitle{第十五章：封装、OCP、DIP 与 LSP}

\question{信息隐藏与封装的含义}
\begin{answer}
信息隐藏关注“隐藏设计决策和变化点”；封装关注“把数据和操作组织在模块或对象内部，并通过接口访问”。封装是实现信息隐藏的重要手段，但有封装语法不代表真正完成信息隐藏；如果暴露内部结构或让外部依赖实现细节，仍然信息隐藏不足。
\end{answer}

\question{OCP 和 DIP}
\begin{answer}
OCP：对扩展开放，对修改关闭。新增业务规则时通过新增类或配置扩展，不频繁修改稳定核心代码。

DIP：高层模块不依赖低层模块，二者都依赖抽象；抽象不依赖细节，细节依赖抽象。典型写法是 Controller 依赖 Service 接口，ServiceImpl 依赖 Repository 接口。
\end{answer}

\question{替换原则例题：MyStack 继承 Vector 是否合理？Employee 继承 Person 是否合理？}
\begin{answer}
\textbf{MyStack extends Vector 不合理。} 栈只允许压栈、弹栈、取大小和判空，但继承 Vector 后会暴露 \texttt{insertElementAt}、\texttt{removeElementAt}、\texttt{elementAt} 等任意位置操作，外部能绕过栈的后进先出约束。这违反 LSP，也违反封装和组合优于继承。修改方式是组合：
\begin{lstlisting}[language=Java]
public class MyStack {
    private final Vector<Object> elements = new Vector<>();

    public void push(Object element) {
        elements.insertElementAt(element, 0);
    }

    public Object pop() {
        Object result = elements.firstElement();
        elements.removeElementAt(0);
        return result;
    }

    public int size() {
        return elements.size();
    }

    public boolean isEmpty() {
        return elements.isEmpty();
    }
}
\end{lstlisting}

\textbf{Employee extends Person 合理。} 只要业务语义满足“每个 Employee 都是 Person”，且 Employee 不改变 Person 的基本契约，就符合 LSP。Employee 可以增加员工号、部门、薪资等属性，但不能让 \texttt{getName()} 的含义变化，也不能增强父类方法的前置条件或削弱后置条件。
\end{answer}

\chaptertitle{第十六章：设计模式}

\question{如何实现可修改性、可扩展性和灵活性？}
\begin{answer}
方法：
\begin{itemize}
  \item 找变化点：算法、产品族、创建逻辑、遍历方式、外部接口。
  \item 抽象稳定接口：让高层代码依赖接口。
  \item 封装变化：把变化逻辑放入独立实现类。
  \item 使用组合和委托：运行时替换实现。
  \item 用测试保护变化：为接口契约和核心场景写测试。
\end{itemize}
\end{answer}

\question{策略模式}
\begin{answer}
定义一组算法，把每个算法封装起来，并让它们可以互换，使算法变化独立于使用算法的客户。

适用场景：同一业务存在多种算法或规则，如会员折扣、支付方式、排序算法、积分规则。

结构：\texttt{Context} 依赖 \texttt{Strategy} 接口；每个具体算法实现一个 \texttt{ConcreteStrategy}。新增算法时新增实现类，不修改 Context。
\end{answer}

\question{抽象工厂模式}
\begin{answer}
提供一个创建一组相关或相互依赖对象的接口，而不指定它们的具体类。

适用场景：需要切换一组产品族，如 MySQL DAO 与 Oracle DAO、Windows UI 控件与 macOS UI 控件。抽象工厂保证同一产品族的一致性。
\end{answer}

\question{单件模式}
\begin{answer}
单件模式保证一个类只有一个实例，并提供全局访问点。适用于全局唯一且无可变共享风险的对象，如只读配置、日志器入口。

风险：容易形成全局状态、隐式耦合和测试困难。课程答题中要说明边界：不要把数据库连接、业务服务和可变缓存随意做成单件。
\end{answer}

\question{迭代器模式}
\begin{answer}
迭代器模式在不暴露聚合对象内部表示的前提下，顺序访问其中元素。适用于隐藏集合内部结构，如数组、链表、分页查询、数据库游标。客户端只依赖 Iterator 接口，不关心底层存储。
\end{answer}

\question{给定场景选择设计模式并写代码，给出代码要求用设计模式改写}
\begin{answer}
答题模板：
\begin{enumerate}
  \item 找变化点：多种算法用 Strategy；一组相关对象创建用 Abstract Factory；全局唯一对象用 Singleton；隐藏遍历结构用 Iterator。
  \item 写出角色：接口、具体实现、上下文或工厂。
  \item 说明原则：OCP、DIP、低耦合、封装变化。
  \item 写测试：新增实现类后旧流程不变。
\end{enumerate}

策略模式骨架：
\begin{lstlisting}[language=Java]
public interface DiscountStrategy {
    Money discount(Order order);
}

public class VipDiscountStrategy implements DiscountStrategy {
    public Money discount(Order order) {
        return order.total().multiply(0.8);
    }
}

public class OrderService {
    private final DiscountStrategy strategy;

    public OrderService(DiscountStrategy strategy) {
        this.strategy = strategy;
    }

    public Money calculatePayable(Order order) {
        return strategy.discount(order);
    }
}
\end{lstlisting}
\end{answer}

\chaptertitle{第十七、十八章：软件构造与代码设计}

\question{构造包含哪些活动？}
\begin{answer}
构造不只是编码，包含：
\begin{itemize}
  \item 编码和单元实现。
  \item 单元测试、调试和代码审查。
  \item 集成、构建和部署制品生成。
  \item 静态分析、持续集成和质量门禁。
  \item 重构、缺陷修复和开发者文档。
\end{itemize}
\end{answer}

\question{名词解释：重构、测试驱动开发、结对编程}
\begin{answer}
\begin{itemize}
  \item 重构：在不改变软件外部行为的前提下改善内部结构。
  \item 测试驱动开发 TDD：Red、Green、Refactor。先写失败测试，再写最少代码通过测试，最后在测试保护下重构。
  \item 结对编程：两名开发者共同完成同一任务，一人驾驶编码，一人观察审查，定期交换角色。
\end{itemize}
\end{answer}

\question{给定代码段，如何发现问题并改进？}
\begin{answer}
从以下角度分析：
\begin{itemize}
  \item 命名：变量、方法和类名是否表达意图。
  \item 简洁性：是否存在长方法、重复代码、过深嵌套。
  \item 职责：是否一个方法承担校验、计算、持久化和展示多个职责。
  \item 控制结构：是否能用早返回、卫语句、多态或表驱动简化。
  \item 数据结构：是否用合适结构替代复杂条件判断。
  \item 异常处理：是否吞掉异常、返回魔法值或缺少错误恢复。
  \item 测试：是否能为正常、边界和异常场景编写测试。
\end{itemize}
\end{answer}

\question{变量使用、语句处理和 How to write unmaintainable code}
\begin{answer}
变量使用：
\begin{itemize}
  \item 缩小作用域，减少可变状态。
  \item 使用有意义命名，避免魔法数字。
  \item 用常量表达业务阈值。
\end{itemize}

语句处理：
\begin{itemize}
  \item 使用卫语句减少嵌套。
  \item 复杂条件提取为命名方法。
  \item 避免超长方法和重复分支。
\end{itemize}

不可维护代码的典型问题：命名无意义、全局变量、重复代码、隐藏副作用、过深嵌套、类型码分支、异常被吞掉、缺少测试。
\end{answer}

\question{防御与错误处理、契约式设计、防御式编程}
\begin{answer}
契约式设计强调前置条件、后置条件和不变量。调用者满足前置条件，被调用者保证后置条件。

防御式编程强调即使输入不可信也要检查和处理：
\begin{itemize}
  \item 校验空值、范围、格式和状态。
  \item 对非法输入给出明确异常或错误结果。
  \item 外部资源使用超时、重试和降级。
  \item 不吞异常，记录必要上下文。
\end{itemize}
\end{answer}

\question{表驱动}
\begin{answer}
表驱动用数据表替代复杂分支，适合规则稳定、分支很多的场景。数据保存的重点是节省计算、减少重复判断并提高可维护性。

月份天数示例：
\begin{lstlisting}[language=Java]
private static final int[] DAYS = {
    0, 31, 28, 31, 30, 31, 30,
    31, 31, 30, 31, 30, 31
};

public int daysOfMonth(int month) {
    if (month < 1 || month > 12) {
        throw new IllegalArgumentException("month");
    }
    return DAYS[month];
}
\end{lstlisting}
\end{answer}

\question{单元测试用例设计}
\begin{answer}
单元测试关注单个类或方法：
\begin{itemize}
  \item 正常输入。
  \item 边界输入。
  \item 非法输入。
  \item 状态变化。
  \item 外部依赖用 Mock 替代。
\end{itemize}

表格格式：编号、前置条件、输入、覆盖点、预期输出。
\end{answer}

\chaptertitle{第十九章：测试}

\question{黑盒测试和白盒测试的常见方法，并比较优缺点}
\begin{answer}
黑盒测试基于规格，不看内部代码。方法包括等价类划分、边界值分析、决策表、状态转换、场景测试。优点是贴近用户需求；缺点是无法直接保证代码路径覆盖。

白盒测试基于代码结构。方法包括语句覆盖、分支覆盖、条件覆盖、路径覆盖。优点是能发现内部逻辑遗漏；缺点是依赖代码结构，不能替代需求层面的功能验证。
\end{answer}

\question{解释并区别语句覆盖、分支覆盖和路径覆盖}
\begin{answer}
\begin{itemize}
  \item 语句覆盖：每条语句至少执行一次。强度最低，不能保证每个判定结果都执行。
  \item 分支覆盖：每个判定的真分支和假分支至少执行一次。强于语句覆盖。
  \item 路径覆盖：覆盖程序中的执行路径。强度高，但路径数量随条件组合快速增长，成本最高。
\end{itemize}
\end{answer}

\question{给定场景，如何选择测试方法并写测试用例？}
\begin{answer}
选择规则：
\begin{itemize}
  \item 给出功能需求：写功能测试，用等价类、边界值、决策表和场景测试。
  \item 给出设计图：写集成测试，关注接口、消息顺序、Stub 和 Driver。
  \item 给出方法代码：写单元测试，关注白盒覆盖、边界和异常。
  \item 给出外部依赖：使用 MockObject 隔离数据库、网络或第三方服务。
\end{itemize}

测试用例格式：
\begin{longtable}{p{0.13\linewidth}p{0.28\linewidth}p{0.25\linewidth}p{0.24\linewidth}}
\toprule
编号 & 输入/前置条件 & 覆盖点 & 预期结果 \\
\midrule
TC-01 & 有效输入 & 正常路径 & 返回成功结果。\\
TC-02 & 最小边界值 & 边界值 & 接受或返回规定错误。\\
TC-03 & 刚好越界 & 无效等价类 & 拒绝并提示原因。\\
TC-04 & 依赖服务失败 & 异常路径 & 回滚或返回失败结果。\\
\bottomrule
\end{longtable}
\end{answer}

\question{JUnit 基本使用方法}
\begin{answer}
JUnit 测试通常包含 Arrange、Act、Assert 三步：
\begin{lstlisting}[language=Java]
@Test
void shouldRejectZeroAmount() {
    RechargeService service = new RechargeService(fakeRepository);

    assertThrows(IllegalArgumentException.class, () -> {
        service.recharge("card-001", BigDecimal.ZERO);
    });
}
\end{lstlisting}
要点：测试名表达行为；每个测试验证一个重点；使用断言检查返回值、异常或状态；外部依赖用 Fake、Stub 或 Mock 替代。
\end{answer}

\chaptertitle{第二十、二十一章：维护与演化}

\question{如何理解软件维护的重要性？}
\begin{answer}
软件交付后仍会因需求变化、运行环境变化、缺陷修复、安全漏洞、性能压力和业务扩展而持续修改。维护阶段通常占据软件生命周期的大量成本。可维护性差会导致修改风险高、缺陷回归多、交付速度慢。
\end{answer}

\question{开发可维护软件的方法}
\begin{answer}
方法：
\begin{itemize}
  \item 清晰需求和可追踪关系。
  \item 分层架构、低耦合、高内聚。
  \item 信息隐藏和稳定接口。
  \item 编码规范、清晰命名和小方法。
  \item 自动化测试、持续集成和回归测试。
  \item 配置管理、版本控制和变更记录。
  \item 必要文档，包括用户文档和系统文档。
\end{itemize}
\end{answer}

\question{演化式生命周期模型}
\begin{answer}
演化式生命周期模型先交付初始版本或原型，再根据用户反馈、环境变化和新需求持续修改。它适合需求不稳定、问题域需要探索的系统。

优点：早反馈，能吸收真实用户经验。缺点：范围控制难，架构规划不足时会退化为不断打补丁。
\end{answer}

\question{用户文档和系统文档}
\begin{answer}
\begin{itemize}
  \item 用户文档：面向使用者，说明安装、登录、功能操作、常见问题和错误处理。
  \item 系统文档：面向开发和维护人员，说明需求、架构、设计、接口、数据库、部署、测试和运维。
\end{itemize}
用户文档帮助正确使用系统，系统文档帮助理解、维护和演化系统。
\end{answer}

\question{逆向工程和再工程}
\begin{answer}
逆向工程：从已有代码或系统中恢复设计、需求、架构或业务规则。

再工程：在理解旧系统的基础上重构、迁移、重写或改造，使系统更易维护、适配新环境或满足新需求。

区别：逆向工程侧重理解和恢复；再工程侧重改造和提升。
\end{answer}

\chaptertitle{第二十二、二十三章：生命周期模型与知识体系}

\question{软件生命周期模型}
\begin{answer}
软件生命周期是软件从提出、开发、运行、维护到退役的完整过程。生命周期模型描述这些活动的组织方式。常见模型包括 Build-and-fix、瀑布模型、增量迭代模型、演化模型、原型模型和螺旋模型。
\end{answer}

\question{解释并比较不同过程模型}
\begin{answer}
\begin{longtable}{p{0.18\linewidth}p{0.31\linewidth}p{0.21\linewidth}p{0.2\linewidth}}
\toprule
模型 & 特征 & 优点 & 缺点 \\
\midrule
Build-and-fix & 直接编码，边做边改。 & 启动快。 & 过程不可控，质量和维护风险高。\\
瀑布模型 & 阶段顺序清楚，文档驱动。 & 适合需求稳定、合同明确项目。 & 反馈晚，早期错误传递到后期。\\
增量模型 & 分批交付可运行功能。 & 早交付，风险分散。 & 需要总体架构和增量规划。\\
迭代模型 & 多轮重复分析、设计、实现和测试。 & 持续改进。 & 迭代控制不当会反复返工。\\
演化模型 & 初始版本随反馈持续演化。 & 适合需求变化。 & 范围和进度控制难。\\
原型模型 & 用原型澄清需求。 & 帮助用户表达需求。 & 抛弃式原型不能直接当最终产品。\\
螺旋模型 & 风险驱动，每轮解决主要风险。 & 适合大型高风险项目。 & 风险分析能力要求高，管理成本高。\\
\bottomrule
\end{longtable}
\end{answer}

\question{给定场景，判定适用的开发过程模型}
\begin{answer}
判定规则：
\begin{itemize}
  \item 需求稳定、合同边界清楚：瀑布模型。
  \item 核心需求清楚、功能可按优先级分批交付：增量迭代模型。
  \item 需求变化频繁、需要边用边改：演化模型。
  \item 用户说不清界面、流程或业务规则：原型模型。
  \item 大型、高风险、技术和需求不确定性高：螺旋模型。
\end{itemize}
\end{answer}

\question{软件工程知识体系的知识域}
\begin{answer}
SWEBOK 覆盖需求、设计、构造、测试、维护、配置管理、工程管理、过程、模型与方法、质量、职业实践、经济学、计算基础、数学基础、工程基础等知识域。

SWEBOK V4 增强关注软件体系结构、软件工程运维、软件安全、AI for SE 和 SE for AI。答题收束句：AI 可以提高需求整理、设计草图、代码生成和审查效率，但工程纪律仍由需求验证、接口约束、测试、配置管理和人工评审保证。
\end{answer}

\chaptertitle{题型和数量}

\question{考试题型和数量如何理解？}
\begin{answer}
原提纲最后写明数量为 8 到 10 道大题，题型为中文试题，全部是大题。复习时按大题组织答案：定义题写关键词和展开句；建模题写过程和模型元素；代码题写违反原则、影响、修改和测试；案例题写问题、原因、改进和验证。
\end{answer}

\end{document}
